\chapter{Communication Complexity — Private Coin Approximation}
%%% generated by the blueprint pipeline from TCSlib.CommunicationComplexity.NewmanTheorem.PrivateCoinApproximation
\section{Overview}

This module develops the coin-approximation machinery used in Newman's theorem.
Starting from a CDF-inversion construction that approximates an arbitrary probability
mass function on $\mathrm{Fin}\,m$ by a uniform distribution over a finite coin tape,
it builds up to the key result that any finite-message private-coin protocol over
arbitrary finite probability spaces can be replaced by one whose randomness comes from
a standard coin tape, at the cost of an additive $\delta$ increase in error.

%%%%%%%%%%%%%%%%%%%%%%%%%%%%%%%%%%%%%%%%%%%%%%%%%%%%%%%%%%%%%%%%%%%%%%%%%%%%%%%
\section{Declarations}

\begin{definition}[Cumulative distribution function on $\mathrm{Fin}\,m$]
\lean{CommunicationComplexity.Internal.cdf}
\label{CommunicationComplexity.Internal.cdf}
\leanok
Given a probability mass function $p$ on $\mathrm{Fin}\,m$ and a natural number $n$,
$\mathrm{cdf}(p, n)$ is the cumulative probability $\sum_{j < n} p(j)$, computed as an
extended non-negative real ($\mathbb{R}_{\ge 0}^\infty$).
\end{definition}

\begin{lemma}[Cumulative distribution function vanishes at zero]
\lean{CommunicationComplexity.Internal.cdf_zero}
\label{CommunicationComplexity.Internal.cdf_zero}
\leanok
\uses{CommunicationComplexity.Internal.cdf}
\difficulty{1}
Let $p$ be a probability mass function on $\mathrm{Fin}\,m$. Then the cumulative
probability $\sum_{j < 0} p(j)$ vanishes; that is, $\mathrm{cdf}(p, 0) = 0$ in
$\bbr_{\ge 0}^\infty$.
\end{lemma}

\begin{lemma}[CDF successor recurrence]
\lean{CommunicationComplexity.Internal.cdf_succ}
\label{CommunicationComplexity.Internal.cdf_succ}
\leanok
\uses{CommunicationComplexity.Internal.cdf}
\difficulty{3}
Let $p$ be a probability mass function on $\mathrm{Fin}\,m$, and let $n \in
\mathrm{Fin}\,m$. Then the cumulative distribution function satisfies $\mathrm{cdf}(p,
n+1) = \mathrm{cdf}(p, n) + p(n)$, where the sums are taken as extended non-negative
reals in $\bbr_{\ge 0}^\infty$.
\end{lemma}

\begin{lemma}[Total cumulative probability equals one]
\lean{CommunicationComplexity.Internal.cdf_one}
\label{CommunicationComplexity.Internal.cdf_one}
\leanok
\uses{CommunicationComplexity.Internal.cdf}
\difficulty{2}
Let $p$ be a probability mass function on the finite set $\mathrm{Fin}\,m = \{0, 1,
\dots, m-1\}$. Then its cumulative distribution function evaluated at $m$ accounts for
the entire mass, $\mathrm{cdf}(p, m) = 1$, this value being computed in the extended
non-negative reals $\bbr_{\ge 0}^\infty$.
\end{lemma}

\begin{lemma}[Monotonicity of the cumulative distribution function]
\lean{CommunicationComplexity.Internal.cdf_mono}
\label{CommunicationComplexity.Internal.cdf_mono}
\leanok
\uses{CommunicationComplexity.Internal.cdf}
\difficulty{3}
Let $p$ be a probability mass function on $\mathrm{Fin}\,m$, and for each natural number
$n$ let $\mathrm{cdf}(p, n) = \sum_{j < n} p(j) \in \bbr_{\ge 0}^\infty$ be the total
mass assigned to the indices below $n$. Then the map $n \mapsto \mathrm{cdf}(p, n)$ is
monotone: whenever $n_1 \le n_2$, one has $\mathrm{cdf}(p, n_1) \le \mathrm{cdf}(p,
n_2)$.
\end{lemma}

\begin{definition}[Inverse CDF]
\lean{CommunicationComplexity.Internal.invCdf}
\label{CommunicationComplexity.Internal.invCdf}
\leanok
\uses{CommunicationComplexity.Internal.cdf}
Given a PMF $p$ on $\mathrm{Fin}\,m$ (with $m \ge 1$) and a value $x \in \mathbb{R}_{\ge 0}^\infty$,
$\mathrm{invCdf}(p, x)$ is the largest index $i \in \mathrm{Fin}\,m$ satisfying
$\mathrm{cdf}(p, i) \le x$, i.e.\ the generalized inverse (quantile) of the CDF.
\end{definition}

\begin{theorem}[Characterization of the inverse CDF]
\lean{CommunicationComplexity.Internal.invCdf_eq_iff}
\label{CommunicationComplexity.Internal.invCdf_eq_iff}
\leanok
\uses{CommunicationComplexity.Internal.cdf, CommunicationComplexity.Internal.cdf_mono, CommunicationComplexity.Internal.cdf_one, CommunicationComplexity.Internal.invCdf}
\difficulty{5}
Let $p$ be a probability mass function on $\mathrm{Fin}\,m = \{0, 1, \dots, m-1\}$ with
$m \ge 1$, and write $F(n) = \sum_{j < n} p(j) \in \bbr_{\ge 0}^{\infty}$ for its
cumulative distribution function. Let $x \in \bbr_{\ge 0}^{\infty}$ satisfy $x < 1$, and
let $i \in \mathrm{Fin}\,m$. Then the largest index whose cumulative probability does
not exceed $x$ equals $i$ if and only if
\[
F(i) \le x < F(i+1).
\]
\end{theorem}

\begin{definition}[Uniform approximation via inverse CDF]
\lean{CommunicationComplexity.Internal.uniformApprox}
\label{CommunicationComplexity.Internal.uniformApprox}
\leanok
\uses{CommunicationComplexity.Internal.invCdf}
Given a PMF $p$ on $\mathrm{Fin}\,m$ and a positive integer $n$, the map
$\mathrm{uniformApprox}(p, n) : \mathrm{Fin}\,n \to \mathrm{Fin}\,m$ sends
$j \mapsto \mathrm{invCdf}(p, j/n)$, discretizing the unit interval into $n$ equally
spaced points and mapping each through the quantile function of $p$.
\end{definition}

\begin{lemma}[Counting bound in a real interval]
\lean{CommunicationComplexity.Internal.card_nat_in_Ico}
\label{CommunicationComplexity.Internal.card_nat_in_Ico}
\leanok
\difficulty{5}
Fix $n \in \bbn$ and real numbers $a \le b$. Among the integers $j \in \{0, 1, \dots,
n-1\}$, the number of those satisfying $a \le j < b$ is at most $b - a + 1$.
\end{lemma}

\begin{theorem}[Uniform approximation error bound]
\lean{CommunicationComplexity.Internal.uniformApprox_approx}
\label{CommunicationComplexity.Internal.uniformApprox_approx}
\leanok
\uses{CommunicationComplexity.Internal.card_nat_in_Ico, CommunicationComplexity.Internal.cdf, CommunicationComplexity.Internal.cdf_mono, CommunicationComplexity.Internal.cdf_one, CommunicationComplexity.Internal.cdf_succ, CommunicationComplexity.Internal.invCdf_eq_iff, CommunicationComplexity.Internal.uniformApprox}
\difficulty{5}
Let $p$ be a probability mass function on $\mathrm{Fin}\,m$ with $m \ge 1$, and let $n
\ge 1$ be an integer. For each $j \in \mathrm{Fin}\,n$, let $\varphi(j) \in
\mathrm{Fin}\,m$ be the largest index $k$ whose cumulative probability $\sum_{\ell < k}
p(\ell)$ does not exceed $j/n$; equivalently, $\varphi$ evaluates the generalized
inverse (quantile function) of the cumulative distribution function of $p$ at the $n$
equally spaced points $0, 1/n, \dots, (n-1)/n$. Then for every index $i \in
\mathrm{Fin}\,m$, the fraction of these $n$ points that $\varphi$ sends to $i$ exceeds
the true probability $p(i)$ by at most $1/n$:
\[
  \frac{\abs{\{\, j \in \mathrm{Fin}\,n : \varphi(j) = i \,\}}}{n}
  \;\le\; p(i) + \frac{1}{n},
\]
where $p(i)$ is read as a real number.
\end{theorem}

\begin{theorem}[Approximating a finite probability space by fair coin flips]
\lean{CommunicationComplexity.Internal.single_coin_approx}
\label{CommunicationComplexity.Internal.single_coin_approx}
\leanok
\uses{CommunicationComplexity.CoinTape, CommunicationComplexity.FiniteProbabilitySpace, CommunicationComplexity.FiniteProbabilitySpace.hasSum_measure_singletons, CommunicationComplexity.FiniteProbabilitySpace.measureReal_preimage_finset, CommunicationComplexity.Internal.uniformApprox, CommunicationComplexity.Internal.uniformApprox_approx, CommunicationComplexity.uniformOn_univ_measureReal_eq_card_filter}
\difficulty{6}
Let $\Omega$ be a finite probability space, and let $\delta > 0$. Then there exist a
natural number $n$ and a map $\varphi : \mathrm{CoinTape}(n) \to \Omega$ from the
length-$n$ binary strings such that, writing $\mu_n$ for the uniform probability measure
on $\mathrm{CoinTape}(n)$ (each of the $2^n$ strings equally likely) and $\mu$ for the
probability measure on $\Omega$, every subset $S \subseteq \Omega$ satisfies
\[
  \mu_n\bigl(\varphi^{-1}(S)\bigr) \;\le\; \mu(S) + \delta.
\]
\end{theorem}

\begin{lemma}[Weighted sum approximation]
\lean{CommunicationComplexity.Internal.weighted_sum_approx}
\label{CommunicationComplexity.Internal.weighted_sum_approx}
\leanok
\difficulty{5}
Let $\alpha$ be a finite type, let $p, q, g : \alpha \to \bbr$, and suppose $0 \le g(a)
\le 1$ for every $a \in \alpha$. Fix a real number $\delta$, and suppose that for every
subset $T \subseteq \alpha$ one has $\sum_{a \in T} p(a) \le \sum_{a \in T} q(a) +
\delta$. Then $\sum_{a} p(a)\,g(a) \le \sum_{a} q(a)\,g(a) + \delta$.
\end{lemma}

\begin{theorem}[Product-space coin approximation]
\lean{CommunicationComplexity.Internal.product_coin_approx}
\label{CommunicationComplexity.Internal.product_coin_approx}
\leanok
\uses{CommunicationComplexity.CoinTape, CommunicationComplexity.FiniteProbabilitySpace, CommunicationComplexity.FiniteProbabilitySpace.measureReal_finset, CommunicationComplexity.FiniteProbabilitySpace.measureReal_iUnion_fintype, CommunicationComplexity.FiniteProbabilitySpace.measureReal_preimage_finset, CommunicationComplexity.FiniteProbabilitySpace.measureReal_prod, CommunicationComplexity.Internal.single_coin_approx, CommunicationComplexity.Internal.weighted_sum_approx}
\difficulty{6}
Let $\Omega_X$ and $\Omega_Y$ be finite probability spaces, and let $\delta > 0$. Give
each coin-tape space $\mathrm{CoinTape}(n)$ the uniform probability measure arising from
$n$ independent fair coin flips (mass $2^{-n}$ on each string), and give every product
space — both $\Omega_X \times \Omega_Y$ and $\mathrm{CoinTape}(n_X) \times
\mathrm{CoinTape}(n_Y)$ — the corresponding product measure. Then there exist $n_X, n_Y
\in \bbn$ and maps $\varphi_X : \mathrm{CoinTape}(n_X) \to \Omega_X$ and $\varphi_Y :
\mathrm{CoinTape}(n_Y) \to \Omega_Y$ such that, writing $\varphi_X \times \varphi_Y$ for
the map $(u,v) \mapsto (\varphi_X(u), \varphi_Y(v))$, every subset $S \subseteq \Omega_X
\times \Omega_Y$ satisfies
\[
  \Pr\!\big[(\varphi_X \times \varphi_Y)^{-1}(S)\big] \;\le\; \Pr[S] + \delta .
\]
\end{theorem}

\begin{definition}[Protocol conversion to coin tape]
\lean{CommunicationComplexity.PrivateCoin.FiniteMessage.Protocol.toCoinTape}
\label{CommunicationComplexity.PrivateCoin.FiniteMessage.Protocol.toCoinTape}
\leanok
\uses{CommunicationComplexity.CoinTape, CommunicationComplexity.Deterministic.FiniteMessage.Protocol.comap, CommunicationComplexity.FiniteProbabilitySpace, CommunicationComplexity.Internal.product_coin_approx, CommunicationComplexity.PrivateCoin.FiniteMessage.Protocol}
Given a finite-message private-coin protocol $p$ over finite probability spaces $\Omega_X$
and $\Omega_Y$, and a tolerance $\delta > 0$, $\mathrm{toCoinTape}(p, \delta)$ produces
$n_X$, $n_Y$, and a new finite-message protocol over $\mathrm{CoinTape}(n_X) \times
\mathrm{CoinTape}(n_Y)$ obtained by precomposing $p$ with the inverse-CDF approximating
maps $\varphi_X$ and $\varphi_Y$.
\end{definition}

\begin{theorem}[Coin-tape conversion preserves communication complexity]
\lean{CommunicationComplexity.PrivateCoin.FiniteMessage.Protocol.toCoinTape_complexity}
\label{CommunicationComplexity.PrivateCoin.FiniteMessage.Protocol.toCoinTape_complexity}
\leanok
\uses{CommunicationComplexity.Deterministic.FiniteMessage.Protocol.complexity, CommunicationComplexity.FiniteProbabilitySpace, CommunicationComplexity.PrivateCoin.FiniteMessage.Protocol, CommunicationComplexity.PrivateCoin.FiniteMessage.Protocol.toCoinTape}
\difficulty{1}
Let $\Omega_X$ and $\Omega_Y$ be finite probability spaces, let $X$, $Y$ be input types
and $\alpha$ an output type, and let $p$ be a private-coin finite-message protocol over
$\Omega_X$, $\Omega_Y$ (that is, a deterministic finite-message protocol on the input
types $\Omega_X \times X$ and $\Omega_Y \times Y$). For every real tolerance $\delta >
0$, the coin-tape protocol produced from $p$ and $\delta$ — the finite-message protocol
over $\mathrm{CoinTape}(n_X) \times X$ and $\mathrm{CoinTape}(n_Y) \times Y$ obtained by
precomposing $p$ with the inverse-CDF approximating maps — has communication complexity
exactly equal to that of $p$.
\end{theorem}

\begin{theorem}[Coin-tape conversion preserves approximate correctness]
\lean{CommunicationComplexity.PrivateCoin.FiniteMessage.Protocol.toCoinTape_approxSatisfies}
\label{CommunicationComplexity.PrivateCoin.FiniteMessage.Protocol.toCoinTape_approxSatisfies}
\leanok
\uses{CommunicationComplexity.CoinTape, CommunicationComplexity.Deterministic.FiniteMessage.Protocol.comap, CommunicationComplexity.Deterministic.FiniteMessage.Protocol.comap_run, CommunicationComplexity.FiniteProbabilitySpace, CommunicationComplexity.Internal.product_coin_approx, CommunicationComplexity.PrivateCoin.FiniteMessage.Protocol, CommunicationComplexity.PrivateCoin.FiniteMessage.Protocol.ApproxSatisfies, CommunicationComplexity.PrivateCoin.FiniteMessage.Protocol.rrun, CommunicationComplexity.PrivateCoin.FiniteMessage.Protocol.toCoinTape}
\difficulty{4}
Let $p$ be a private-coin finite-message protocol over input types $X$, $Y$ and output
type $\alpha$, whose private randomness for Alice and Bob is drawn from finite
probability spaces $\Omega_X$ and $\Omega_Y$. Let $Q : X \to Y \to \alpha \to
\mathrm{Prop}$ be a predicate on input–output triples, and let $\varepsilon$ and
$\delta$ be real numbers with $\delta > 0$. Suppose $p$ $\varepsilon$-satisfies $Q$;
that is, for every input pair $(x, y)$ the probability, over Alice's and Bob's private
coins $(\omega_x, \omega_y)$, that the output of $p$ fails to satisfy $Q\,x\,y$ is at
most $\varepsilon$. Then the protocol obtained by converting $p$ at tolerance $\delta$ —
replacing the private randomness by fair coin tapes of some lengths $n_X$ and $n_Y$ —
$(\varepsilon + \delta)$-satisfies $Q$.
\end{theorem}
